% week10_event_driven.tex -- CS2 Week 10 Monday, Oct 19 2026.
% Event-driven programming: Tkinter and a web page over the same Inventory model.
% Code is copied from week_10/compact_gui/ (inventory.py, test_inventory.py, view.py, view_web.py).
\documentclass[aspectratio=169]{beamer}
\usepackage{../theme/cs2theme}
\usepackage{../theme/cs2commands}

\setcsweek{Week 10}

\title{Event-Driven Programming}
\subtitle{Tkinter and a web page, same logic}
\author{Computer Science II \textbullet\ Dr. Jeremy P. Evert}
\date{October 19, 2026}

\begin{document}

\begin{frame}[plain]
  \titlepage
  \note{Headline today is event-driven programming, not Tkinter. Tkinter is one vehicle, a web page is another. The logic underneath is the same, and Python setup trouble is not the lesson.}
\end{frame}

\begin{destinationframe}{Today's destination}
  \begin{itemize}
    \item Map an event: source, handler, state, output.
    \item Trace it from the event to the refreshed display.
    \item Test the rule with no UI.
    \item Verify the handler only routes.
    \item Swap the front end without touching the model.
  \end{itemize}
  \instructorcue{Say it out loud: the front end is a vehicle. The model is what we keep.}
  \note{If time runs short, cut the dependencies slide, never the trace. The gate is the Week 10 checkpoint; this deck is the pointer to the idea, not the rubric.}
\end{destinationframe}

\begin{conceptframe}{What is an event-driven program?}
  \begin{description}
    \item[Cycle] Something happens $\rightarrow$ a handler runs $\rightarrow$ state changes $\rightarrow$ the display refreshes.
    \item[Examples] A key press, a timer, a network packet, a button click, a web request.
  \end{description}
  \instructorcue{Ask: which of these have you already written a handler for?}
  \note{The program does not drive the order; events do. The loop (Tk mainloop, a server loop) waits, and calls your handler when something arrives.}
\end{conceptframe}

\begin{conceptframe}{Cousins beyond GUIs}
  \begin{description}
    \item[Hardware interrupts] A device signals; the CPU runs a handler. That handler is a callback.
    \item[State machines] FPGA and digital designs often react to input signals and change state.
  \end{description}
  \instructorcue{Say ``cousins.'' Same shape of idea, different machinery.}
  \note{Be careful here. Do not say CPUs or FPGAs come from event-driven programming. They are cousins: a signal arrives, a defined response runs, state changes. That resemblance is why the pattern feels familiar in hardware courses.}
\end{conceptframe}

\begin{frame}{Map the event: Inventory restock}
\begin{center}
\begin{tabular}{@{}ll@{}}
\hline
\textbf{Event} & Restock is requested \\
\textbf{Source} & Button click (Tkinter) or POST /restock (web) \\
\textbf{Handler} & \texttt{on\_restock} or \texttt{do\_POST} \\
\textbf{State} & The model's \texttt{items} list \\
\textbf{Output} & ``Items: n'' shown to the user \\
\hline
\end{tabular}
\end{center}
\instructorcue{Have students fill in this table for their own world noun first.}
\note{Five words to map any event. Students reuse this table in the World Bible entry.}
\end{frame}

\begin{frame}[fragile]{The model stays the same}
\lstset{basicstyle=\ttfamily\scriptsize}
\begin{columns}[T]
\begin{column}{0.5\textwidth}
\begin{lstlisting}
# inventory.py
class Inventory:
    def __init__(self):
        self.items = []

    def restock(self, name):
        self.items.append(name)
        return len(self.items)

    def count(self):
        return len(self.items)
\end{lstlisting}
\end{column}
\begin{column}{0.5\textwidth}
\begin{lstlisting}
# test_inventory.py
class InventoryTest(unittest.TestCase):
    def test_restock_adds_one(self):
        inv = Inventory()
        self.assertEqual(inv.restock("rope"), 1)
        self.assertEqual(inv.count(), 1)

    def test_three_restocks(self):
        inv = Inventory()
        for _ in range(3):
            inv.restock("rope")
        self.assertEqual(inv.count(), 3)
\end{lstlisting}
\end{column}
\end{columns}
\note{Run the test live: python3 -m unittest. It passes with no UI, no window, no server. Both test file imports are `import unittest` and `from inventory import Inventory`, trimmed here to fit.}
\end{frame}

\begin{frame}[fragile]{Pass 1: the Tkinter view}
\lstset{basicstyle=\ttfamily\scriptsize}
\predictkicker
\begin{lstlisting}
import tkinter as tk
from inventory import Inventory

model, root = Inventory(), tk.Tk()
text = tk.StringVar(value="Items: 0")
tk.Label(root, textvariable=text).pack()

def on_restock():
    n = model.restock("rope")   # route to the model
    text.set(f"Items: {n}")     # then refresh the view

tk.Button(root, text="Restock", command=on_restock).pack()
root.mainloop()
\end{lstlisting}
  \textbf{Predict:} three clicks on Restock. What does the label say?
  \note{Answer: Items: 3. The number lives in the model, the label only shows it. This is view.py from week_10/compact_gui; the file guards mainloop with if __name__ == "__main__".}
\end{frame}

\begin{frame}[fragile]{Pass 2: a web page view}
\lstset{basicstyle=\ttfamily\tiny}
\begin{lstlisting}
model = Inventory()

PAGE = """<!doctype html><title>Inventory</title>
<p>Items: {n}</p>
<form method="post" action="/restock"><button>Restock</button></form>"""

class Handler(BaseHTTPRequestHandler):
    def do_POST(self):
        if self.path == "/restock":
            n = model.restock("rope")     # route to the model
            print(f"request -> handler -> model.restock -> {n} items")
            self.send_response(303)       # then refresh the view
            self.send_header("Location", "/")
            self.end_headers()
        else:
            self.send_error(404)
\end{lstlisting}
  \textbf{Same Inventory.} Python standard library only (\texttt{http.server}); run \texttt{python3 view\_web.py}, open \texttt{localhost:8000}.
  \note{This is view_web.py from week_10/compact_gui, with do_GET and the page-writing helper left out to fit the slide. Point out: the handler body has the same two jobs as on_restock. Route to the model, then refresh.}
\end{frame}

\begin{frame}{Same logic, two front ends}
\begin{center}
\begin{tabular}{@{}lll@{}}
\hline
 & \textbf{Tkinter} & \textbf{Web page} \\
\hline
\textbf{Event source} & Button click & POST /restock \\
\textbf{Handler wiring} & \texttt{command=on\_restock} & \texttt{do\_POST} \\
\textbf{Display} & \texttt{text.set} on a label & Page shows the count \\
\hline
\textbf{Unchanged} & \multicolumn{2}{l}{\texttt{Inventory}, its test, the restock rule} \\
\hline
\end{tabular}
\end{center}
\instructorcue{Ask: what did we have to edit in inventory.py to add the web page? (Nothing.)}
\note{The top three rows are what changed. The last row is what the whole design is for.}
\end{frame}

\begin{frame}[fragile]{Trace the event in each}
\codeframekicker
\textbf{Tkinter:} click $\rightarrow$ callback $\rightarrow$ model.restock $\rightarrow$ 1 $\rightarrow$ refresh

\vspace{2mm}
\textbf{Web:} request $\rightarrow$ handler $\rightarrow$ model.restock $\rightarrow$ 1 $\rightarrow$ response
\begin{lstlisting}
request -> handler -> model.restock -> 1 items
\end{lstlisting}
\instructorcue{Students paste their own trace into the World Bible as evidence.}
\note{The web line is printed by view_web.py itself. For Tkinter, add a print in on_restock the way last session's trace slide did. After one event, the count is 1 in both. This trace is the evidence artifact for the World Bible.}
\end{frame}

\begin{frame}[fragile]{Debug and verify: break the boundary}
\begin{comparebody}
  \begin{compareleft}{Boundary broken}
\begin{lstlisting}
items = []   # view owns data
def on_restock():
    items.append("rope")
    n = len(items)
    text.set(f"Items: {n}")
\end{lstlisting}
  \end{compareleft}
  \begin{compareright}{Boundary kept}
\begin{lstlisting}
model = Inventory()
def on_restock():
    n = model.restock("rope")
    text.set(f"Items: {n}")
\end{lstlisting}
  \end{compareright}
\end{comparebody}
\instructorcue{If you cannot test a rule without the UI, the rule is in the wrong place.}
\note{Ask how to test the left version without a window. The refactor is moving two lines. A handler should only route: call one model method, then refresh.}
\end{frame}

\begin{frame}{Dependencies are an avenue, not the point}
  \begin{itemize}
    \item Tkinter needs \texttt{python3-tk} on Ubuntu.
    \item The web version needs nothing beyond Python.
    \item Setup trouble is incidental. The model and its test run either way.
  \end{itemize}
  \instructorcue{If a laptop will not open a window, switch to the web pass and keep going.}
  \note{Keep this to one minute. Tell students where to look (check_my_setup.sh in week_10/compact_gui) and move on. Do not turn the class into an installation clinic.}
\end{frame}

\begin{designdecisionframe}{Peer swap and World Bible}
  \textbf{Contract test:} trade model files with a partner, then run \texttt{python3 contract\_test.py their\_module}.
  \begin{itemize}
    \item It prints ``contract holds'': your front end works on their model.
    \item It fails: decide whose method names or return values should change.
  \end{itemize}
  \textbf{World Bible entry:} the rule you moved out of a handler, and your trace.
  \note{contract_test.py is in week_10/compact_gui. It checks restock and count exist, restock returns an int, and count is 1 after one restock. The failure is the lesson. The World Bible entry records the rule moved, the date, and the trace.}
\end{designdecisionframe}

\begin{recapframe}{Recap}
  \begin{itemize}
    \item An event runs a handler, the handler changes state, the display refreshes.
    \item The model and its test do not change when the front end does.
    \item Handlers route; rules live in the model.
    \item Tkinter and a web page are two vehicles for one design.
  \end{itemize}
\end{recapframe}

\begin{exitframe}{Exit question}
  \Large\color{cs2primary}
  \textbf{Which part of your program would you keep if the front end changed tomorrow, and why?}

  \vspace{4mm}
  \normalsize\normalfont\color{cs2ink}
  \textbf{Gate:} due Friday 11:59 PM Central.
\end{exitframe}

\end{document}
